% Autorski schemat pojęciowy. Podstawę stanowią cytowane opisy metod w rozdziałach 2–3.
% Strzałki wskazują wspólną ideę; układ nie przedstawia rankingu ani chronologii.
\begin{minipage}[t]{0.30\linewidth}
\centering
\textbf{Przetwarzanie cech}\par\medskip
\fbox{\parbox{0.85\linewidth}{\centering TabNet\\sekwencyjny wybór cech}}\par\medskip
\fbox{\parbox{0.85\linewidth}{\centering FT-Transformer\\uwaga między tokenami}}
\end{minipage}\hfill
\begin{minipage}[t]{0.30\linewidth}
\centering
\textbf{Sieci gęste}\par\medskip
\fbox{\strut MLP}\par
$\Downarrow$\par
\textit{wspólna podstawa}\par\smallskip
\fbox{\parbox{0.85\linewidth}{\centering
RealMLP\\przygotowanie danych i~procedura uczenia\par\medskip
TabM\\podmodele ze wspólnymi parametrami}}
\end{minipage}\hfill
\begin{minipage}[t]{0.30\linewidth}
\centering
\textbf{Sąsiedztwo}\par\medskip
\fbox{\strut KNN}\par
$\Downarrow$\par
\textit{wspólna idea}\par\smallskip
\fbox{\parbox{0.85\linewidth}{\centering ModernNCA\\sąsiedztwo w~uczonej reprezentacji}}
\end{minipage}
\par\bigskip
\fbox{\parbox{0.92\linewidth}{\centering
\textbf{Wiedza przenoszona między zadaniami}\\
Pretrenowanie $\longrightarrow$ TabPFN-3.5 $\longleftarrow$ oznakowany kontekst zadania}}
\par\medskip
\textbf{Pozostałe odniesienia konwencjonalne}\\
Regresja logistyczna: model liniowy; XGBoost i~CatBoost: wzmacniane drzewa.
